\chapter*{ПРИЛОГ А — ФИКСИРАНИ ПАРАМЕТРИ}
\addcontentsline{toc}{chapter}{Прилог А — фиксирани параметри}
\label{app:parametri}

Табела обухвата све параметре који су фиксирани пре покретања експеримента и који се
не мењају ни у једном услову, осим оних који сами представљају експерименталне факторе.

\begin{center}
\begin{longtable}{L{4.4cm} L{3.4cm} L{5.2cm}}
\caption{Фиксирани параметри експеримента}
\label{tab:fiksirani}\\
\toprule
Параметар & Вредност & Разлог фиксирања \\
\midrule
\endfirsthead
\toprule
Параметар & Вредност & Разлог фиксирања \\
\midrule
\endhead
\midrule
\multicolumn{3}{r}{\footnotesize наставак на следећој страни}\\
\endfoot
\bottomrule
\endlastfoot
Период узорка & 2010–2025 & расположивост једноминутне историје \\
Период евалуације & 2015–2025 & 2014. је година загревања клизних прозора \\
Праг за TrBPV & 3,0 дневне робусне девијације & увођење као четврте димензије дало би неуравнотежен план \\
Праг наивног теста & 3 стандардне девијације, двострано & доња граница квалитета, није озбиљна метода \\
Ниво значајности BNS теста & $\alpha = 0{,}001$, једнострано & строг праг смањује број лажних скокова у дугом узорку \\
Ниво Гумбелове критичне вредности (LM) & 0,999 & контрола грешке унутар дана преко свих интрадневних тестова \\
Стопа лажних открића (LM-FDR) & $q = 0{,}05$ & контролише удео лажних међу проглашеним данима; при овим нивоима блаже од Гумбеловог прага \\
Локални прозор LM теста & 4 сесије на свакој мрежи & да се са мрежом мења само фреквенција, а не и дужина прозора \\
Клизни прозор FDR корекције & 250 дана & приближно једна година трговања; спречава гледање унапред \\
Минимална покривеност дана & 80\% очекиваног броја једноминутних свећица & филтрирање непотпуних сесија, иста одлука на свакој мрежи \\
Апсолутни праг опсервација & 20 приноса на датој мрежи & обезбеђује стабилну оцену реализоване квартичности \\
Поновно оцењивање модела & сваких 21 дан трговања & приближно месечно; улази се ажурирају свакодневно \\
Клизни прозор за корекцију и рекалибрацију & 252 дана ван-узорачких грешака & приближно једна година; искључиво прошле грешке \\
Хиперпараметри стабала & 200 стабала, дубина 4, 15 листова, стопа 0,05 & подешавање би мешало труд у претрази са ефектом архитектуре \\
Контролна тачка претходно обученог модела & \texttt{ibm-granite/\allowbreak granite-\allowbreak timeseries-\allowbreak ttm-r2}, издање \texttt{512-192-r2} & различита издања имају различите скупове предтренирања \\
Тежине хибридног модела & 0,5 и 0,5 & оптимизација би одразила претрагу, не архитектуру \\
Ниво скупа поверења модела & $\alpha = 0{,}10$ & уобичајен ниво у литератури о скупу поверења \\
Ниво Диболд–Маријановог теста & $\alpha = 0{,}05$, уз контролу стопе лажних открића & велики број поређења захтева контролу породице \\
Граница режима & медијана $\RV_m$ у првом периоду обуке, мрежа 5m & фиксна граница чини ознаку режима једнозначном \\
Ширина појаса локалног Витловог естиматора & $m = \lfloor n^{0{,}65}\rfloor$ & уобичајен избор облика $n^{\alpha}$ \\
Број пермутација у плацебо тесту & 100, унутар симбола & расподела разлика под нултом хипотезом \\
Ниво вредности при ризику & 1\% & стандард у регулаторној пракси \\
\end{longtable}
\end{center}
